% GENERATED FILE. Edit canonical lessons, then run npm run generate:books.
% canonical-source-hash: fnv1a64:9c94154cff535b65
% canonical-lessons: GE-C367-eventuell, GE-C367-moeglicherweise, GE-C367-zumindest, GE-C367-immerhin, GE-C367-womoeglich

\chapter{Vielleicht}
\label{ch:ge-vielleicht}
\hlchaptermodality{367}

\begin{chapteropening}
\textbf{By the end of this chapter:} \emph{I can say that something might happen: possibly, maybe, perhaps even, and add at any rate and after all.}

Complete the last lesson of chapter 367: I can say that something might happen: possibly, maybe, perhaps even, and add at any rate and after all.
\end{chapteropening}

\section[eventuell]{eventuell — possibly (not eventually); perhaps}

\label{lesson:GE-C367-eventuell}



\begin{quote}
Before the new one: say the German for if applicable, then the German for at least.
\end{quote}

\subsection*{You'll want to know: eventuell}
\textbf{eventuell} — \textquotedblleft{}possibly (not eventually); perhaps\textquotedblright{}.

\begin{quote}
Eventuell komme ich später, ich weiß es noch nicht.
\end{quote}

\begin{quote}
\emph{I may possibly come later; I do not know yet.}
\end{quote}

A false friend: \emph{eventuell} means possibly, not eventually. Eventually is \emph{irgendwann}.

\begin{cousinweb}
\emph{eventuell} was borrowed from French \emph{éventuel} \textquotedblleft{}possible\textquotedblright{}, a learned formation on Latin \emph{ēventus} \textquotedblleft{}an outcome, an event\textquotedblright{}, from \emph{ēvenīre} \textquotedblleft{}to come out, to happen\textquotedblright{} (Lewis \& Short).
\end{cousinweb}

\begin{grammarlens}[title={What you've built}]
One more word for telling, explaining and describing at B1.
\end{grammarlens}

\subsection*{Your turn}
\begin{itemize}
\raggedright
  \item \emph{Say it:} \emph{eventuell}
  \item \emph{Say it:} \emph{eventuell}, once more
  \item \emph{Say it:} a sentence of your own with \emph{eventuell}
  \item \emph{Recall:} say the German for if applicable, then the German for at least, then say \emph{eventuell} again
\end{itemize}

\begin{tcolorbox}[breakable,colback=teal!4,colframe=teal!35!black,title={Before you move on}]
What does eventuell mean? (\textquotedblleft{}Possibly (not eventually); perhaps\textquotedblright{}.) Say it once more.
\end{tcolorbox}
\section[möglicherweise]{möglicherweise — maybe, possibly}

\label{lesson:GE-C367-moeglicherweise}



\begin{quote}
Before the new one: say the German for at least, then the German for possibly.
\end{quote}

\subsection*{You'll want to know: möglicherweise}
\textbf{möglicherweise} — \textquotedblleft{}maybe, possibly\textquotedblright{}.

\begin{quote}
Möglicherweise ist der Zug schon weg.
\end{quote}

\begin{quote}
\emph{Maybe the train has already gone.}
\end{quote}

Built like \emph{glücklicherweise}: \emph{möglich} + \emph{-erweise}. The verb comes right after it.

\begin{cousinweb}
\emph{möglicherweise} is built from \emph{möglich} \textquotedblleft{}possible\textquotedblright{}, of the family of \emph{mögen} and Proto-Germanic \emph{maganą} \textquotedblleft{}to be able\textquotedblright{} (EtymDB), and \emph{Weise} \textquotedblleft{}way, manner\textquotedblright{}.
\end{cousinweb}

\begin{grammarlens}[title={What you've built}]
One more word for telling, explaining and describing at B1.
\end{grammarlens}

\subsection*{Your turn}
\begin{itemize}
\raggedright
  \item \emph{Say it:} \emph{möglicherweise}
  \item \emph{Say it:} \emph{möglicherweise}, once more
  \item \emph{Say it:} a sentence of your own with \emph{möglicherweise}
  \item \emph{Recall:} say the German for at least, then the German for possibly, then say \emph{möglicherweise} again
\end{itemize}

\begin{tcolorbox}[breakable,colback=teal!4,colframe=teal!35!black,title={Before you move on}]
What does möglicherweise mean? (\textquotedblleft{}Maybe, possibly\textquotedblright{}.) Say it once more.
\end{tcolorbox}
\section[zumindest]{zumindest — at any rate, at least}

\label{lesson:GE-C367-zumindest}



\begin{quote}
Before the new one: say the German for possibly, then the German for maybe.
\end{quote}

\subsection*{You'll want to know: zumindest}
\textbf{zumindest} — \textquotedblleft{}at any rate, at least\textquotedblright{}.

\begin{quote}
Zumindest hat er sich entschuldigt.
\end{quote}

\begin{quote}
\emph{At any rate, he apologized.}
\end{quote}

Close to \emph{wenigstens}. The verb comes right after it: \emph{zumindest hat er …}.

\begin{cousinweb}
\emph{zumindest} grew out of \emph{zum mindesten} \textquotedblleft{}at the least\textquotedblright{}, built on \emph{minder} \textquotedblleft{}less\textquotedblright{}, from Old High German \emph{minniro} and Proto-Germanic \emph{minnizô} (EtymDB).
\end{cousinweb}

\begin{grammarlens}[title={What you've built}]
One more word for telling, explaining and describing at B1.
\end{grammarlens}

\subsection*{Your turn}
\begin{itemize}
\raggedright
  \item \emph{Say it:} \emph{zumindest}
  \item \emph{Say it:} \emph{zumindest}, once more
  \item \emph{Say it:} a sentence of your own with \emph{zumindest}
  \item \emph{Recall:} say the German for possibly, then the German for maybe, then say \emph{zumindest} again
\end{itemize}

\begin{tcolorbox}[breakable,colback=teal!4,colframe=teal!35!black,title={Before you move on}]
What does zumindest mean? (\textquotedblleft{}At any rate, at least\textquotedblright{}.) Say it once more.
\end{tcolorbox}
\section[immerhin]{immerhin — after all, still (on the bright side)}

\label{lesson:GE-C367-immerhin}



\begin{quote}
Before the new one: say the German for maybe, then the German for at any rate.
\end{quote}

\subsection*{You'll want to know: immerhin}
\textbf{immerhin} — \textquotedblleft{}after all, still (on the bright side)\textquotedblright{}.

\begin{quote}
Wir haben verloren. Immerhin haben wir ein Tor geschossen.
\end{quote}

\begin{quote}
\emph{We lost. Still, we scored a goal.}
\end{quote}

It names the good side of a bad result. The verb comes right after it.

\begin{cousinweb}
\emph{immerhin} joins \emph{immer} \textquotedblleft{}always\textquotedblright{}, from Middle High German \emph{immer} (EtymDB), to \emph{hin} \textquotedblleft{}onward, away\textquotedblright{}: whatever happens, this much stands.
\end{cousinweb}

\begin{grammarlens}[title={What you've built}]
One more word for telling, explaining and describing at B1.
\end{grammarlens}

\subsection*{Your turn}
\begin{itemize}
\raggedright
  \item \emph{Say it:} \emph{immerhin}
  \item \emph{Say it:} \emph{immerhin}, once more
  \item \emph{Say it:} a sentence of your own with \emph{immerhin}
  \item \emph{Recall:} say the German for maybe, then the German for at any rate, then say \emph{immerhin} again
\end{itemize}

\begin{tcolorbox}[breakable,colback=teal!4,colframe=teal!35!black,title={Before you move on}]
What does immerhin mean? (\textquotedblleft{}After all, still (on the bright side)\textquotedblright{}.) Say it once more.
\end{tcolorbox}
\section[womöglich]{womöglich — perhaps even, quite possibly}

\label{lesson:GE-C367-womoeglich}



\begin{quote}
Before the new one: say the German for at any rate, then the German for after all.
\end{quote}

\subsection*{You'll want to know: womöglich}
\textbf{womöglich} — \textquotedblleft{}perhaps even, quite possibly\textquotedblright{}.

\begin{quote}
Der Himmel ist ganz dunkel, womöglich gibt es heute noch ein Gewitter.
\end{quote}

\begin{quote}
\emph{The sky is very dark; there may even be a thunderstorm today.}
\end{quote}

Often with a hint of worry: something unwelcome might happen. The verb comes right after it.

\begin{cousinweb}
\emph{womöglich} joins \emph{wo} \textquotedblleft{}where\textquotedblright{} and \emph{möglich} \textquotedblleft{}possible\textquotedblright{}, of the family of \emph{mögen} and Proto-Germanic \emph{maganą} \textquotedblleft{}to be able\textquotedblright{} (EtymDB): \textquotedblleft{}if at all possible\textquotedblright{}, then \textquotedblleft{}perhaps even\textquotedblright{}.
\end{cousinweb}

\begin{grammarlens}[title={What you've built}]
One more word for telling, explaining and describing at B1.
\end{grammarlens}

\subsection*{Your turn}
\begin{itemize}
\raggedright
  \item \emph{Say it:} \emph{womöglich}
  \item \emph{Say it:} \emph{womöglich}, once more
  \item \emph{Say it:} a sentence of your own with \emph{womöglich}
  \item \emph{Recall:} say the German for at any rate, then the German for after all, then say \emph{womöglich} again
\end{itemize}

\begin{tcolorbox}[breakable,colback=teal!4,colframe=teal!35!black,title={Before you move on}]
What does womöglich mean? (\textquotedblleft{}Perhaps even, quite possibly\textquotedblright{}.) Say it once more.
\end{tcolorbox}
